% ============================================================
% GENERAL CONCLUSION
% ============================================================
\chapter*{General Conclusion}
\addcontentsline{toc}{chapter}{General Conclusion}

IGS Marketplace was successfully designed and developed over five sprints, delivering Tunisia's first dedicated 3D digital asset platform. The system provides a six-level Role-Based Access Control system, an interactive in-browser 3D asset viewer, Stripe payment integration, AI-powered review sentiment analysis, a personalized recommendation engine, a floating chatbot assistant powered by Qwen2.5-72B-Instruct, and a companion internal workspace application — all within a cohesive full-stack ecosystem built on React, Node.js/Express, and PostgreSQL.

Several key technical insights emerged throughout the project. Designing the permission system during Sprint 1 and enforcing it consistently across all subsequent sprints was the single most impactful architectural decision — retrofitting access control would have been significantly more costly. The three-tier layered architecture (Routes → Controllers → Repositories) proved highly maintainable for a REST API at this level of complexity, and applying Scrum as a solo developer, with supervisors occupying the Product Owner and Scrum Master roles, effectively structured six months of work into reviewable, deliverable-focused increments.

Looking ahead, several directions are planned for the platform's continued evolution. The internal workspace application will grow into a full \textbf{ERP system} with richer project management, HR, and operational modules, while its customer-facing side will expand into a comprehensive \textbf{CRM} supporting advanced customer segmentation, relationship tracking, and communication workflows. On the payment front, \textbf{Flouci} integration for Tunisian dinar transactions is already partially in place and will be fully activated once the required administrative and regulatory procedures are completed. The asset catalog will be broadened with new content categories to serve a wider range of creative disciplines. A \textbf{bundle system} will allow sellers to group related assets and offer them at a combined price, increasing average order value and discoverability. Finally, the \textbf{voucher system} — currently scoped to internal users — will be extended to all user types, enabling promotional campaigns for members, sellers, and subscribers alike. These evolutions position the IGS Marketplace as a long-term, scalable platform well beyond its current scope.
